\documentclass{article}
\usepackage{amsmath, amssymb, amsthm}
\usepackage{hyperref}
\usepackage{geometry}
\geometry{margin=1in}

\title{Erd\H{o}s Problem \#452}
\date{Last updated: 2025-08-31}
\author{Source: erdosproblems.com}

\begin{document}
\maketitle

\noindent\textbf{Status:} OPEN \quad \textbf{No prize}

\noindent\textbf{Tags:} number theory

\medskip
\noindent\textbf{Problem Statement:}

Let $\omega(n)$ count the number of distinct prime factors of $n$. What is the size of the largest interval $I\subseteq [x,2x]$ such that $\omega(n)>\log\log n$ for all $n\in I$?

\bigskip
\noindent\textbf{Remarks:}

Erd\H{o}s \cite{Er37} proved that the density of integers $n$ with $\omega(n)>\log\log n$ is $1/2$. The Chinese remainder theorem implies that there is such an interval with\[\lvert I\rvert \geq (1+o(1))\frac{\log x}{(\log\log x)^2}.\]It could be true that there is such an interval of length $(\log x)^{k}$ for arbitrarily large $k$.

\bigskip
\noindent\textbf{References:}

[Er37] Erd\H{o}s, Paul, Note on the number of prime divisors of integers. J. London Math. Soc. (1937), 308-314.

\bigskip
\noindent\small{Source: \url{https://www.erdosproblems.com/452}}
\end{document}
